\section{Subsolutions for an individual density}
\label{sec:individual-density-subsolutions}

For one fixed density, the normalizers of its bounded truncations give an exact test for the existence of a positive scaled subsolution. A cluster limit attains the threshold because density lower bounds survive $L^1$ convergence of potentials through upper envelopes.

Let $0\leq f\in L^1(\mathrm{d}V)$ have positive integral, and put
\begin{equation}\label{eq:individual-truncated-normalizer}
 f_k=\min(f,k),\qquad
 b_h(f)=\lim_{k\to\infty}c_h(f_k).
\end{equation}
The limit exists in $[0,+\infty)$ because the constants decrease and $c_h(f_1)$ is finite. The order lemma also gives
\begin{equation}\label{eq:individual-bounded-tests}
 b_h(f)=\inf\{c_h(q):0\leq q\leq f,\ q\in L^\infty,
                                      \ \textstyle\int_Xq\,\mathrm{d}V>0\}.
\end{equation}
In fact every such $q$ is bounded above by $f_k$ for sufficiently large $k$, and the $f_k$ themselves are admissible tests.

A lower bound for a Monge--Amp\`ere density survives $L^1$ convergence of potentials. This uses upper envelopes and damped continuity, rather than continuity of the original Monge--Amp\`ere measures.

\Needspace{8\baselineskip}
\begin{lemma}\label{lem:density-subsolution-lone-closure}
Let $s_j,s\in\PSH(X,h)$ have a common upper bound and satisfy $s_j\to s$ in $L^1(X)$. Suppose $g_j,g\geq0$, $g_j\to g$ in $L^1(\mathrm{d}V)$, and
\[
 \mathrm{MA}^{\mathrm{ext}}_h(s_j)\geq g_j\,\mathrm{d}V.
\]
Then $\mathrm{MA}^{\mathrm{ext}}_h(s)\geq g\,\mathrm{d}V$.
\end{lemma}
\begin{proof}
Translate the potentials to be nonpositive and extract a subsequence for which both convergences hold almost everywhere. Set
\[
 a_j=\inf_{\ell\geq j}g_\ell,\qquad
 w_j=\left(\sup_{\ell\geq j}s_\ell\right)^*.
\]
Then $0\leq a_j\uparrow g$ almost everywhere and $w_j\downarrow s$. For fixed $j$, the finite maxima $\max_{j\leq\ell\leq N}s_\ell$ have measures at least $a_j\,\mathrm{d}V$ by ordered contact. They increase to $w_j$ after regularization, with the genuine lower bound $s_j$ and upper bound zero. The convergence holds in capacity by \cref{lem:monotone-psh-capacity}. Damped continuity therefore gives
\[
 \mathrm{e}^{w_j}\mathrm{MA}^{\mathrm{ext}}_h(w_j)
       \geq \mathrm{e}^{w_j}a_j\,\mathrm{d}V.
\]
Indeed the right sides for the finite maxima converge by dominated convergence, since $a_j\leq g$ almost everywhere. Divide by the positive weight off the pole set to obtain $\mathrm{MA}^{\mathrm{ext}}_h(w_j)\geq a_j\,\mathrm{d}V$.

Now $w_j\downarrow s$ is again convergence in capacity with a genuine lower bound. Apply damped continuity once more. Since $0\leq\mathrm{e}^{w_j}a_j\leq g$, dominated convergence gives
\[
 \mathrm{e}^s\mathrm{MA}^{\mathrm{ext}}_h(s)
       \geq\mathrm{e}^s g\,\mathrm{d}V.
\]
Divide on $\{s>-N\}$ and let $N\to\infty$. Both sides of the desired inequality give zero mass to the pole set. Restore the common translation.
\end{proof}

\begin{theorem}\label{thm:individual-density-subsolution}\index{Subsolution!individual density threshold}
For every such $f$,
\[
 b_h(f)=\max\{a\geq0:\text{there exists }u\in\PSH(X,h)
       \text{ with }\mathrm{MA}^{\mathrm{ext}}_h(u)
                         \geq af\,\mathrm{d}V\}.
\]
In particular, a positive scaled subsolution exists if and only if $b_h(f)>0$.
\end{theorem}
\begin{proof}
If $\mathrm{MA}^{\mathrm{ext}}_h(u)\geq af\,\mathrm{d}V$ with $a>0$, then it dominates $af_k\,\mathrm{d}V$. \Cref{lem:bounded-normalizer-subsolution} gives $a\leq c_h(f_k)$ for every $k$, hence $a\leq b_h(f)$. If $b_h(f)=0$, the zero potential realizes the coefficient zero.

\Needspace{7\baselineskip}
Suppose $b=b_h(f)>0$. Choose continuous normalized solutions
\[
 \mathrm{MA}_h(u_k)=c_h(f_k)f_k\,\mathrm{d}V,
 \qquad \sup_Xu_k=0.
\]
After extraction, $u_{k_j}\to u$ in $L^1$, where $u$ is a genuine normalized $h$-plurisubharmonic function. The densities $c_h(f_{k_j})f_{k_j}$ converge to $bf$ in $L^1$: they are dominated by $c_h(f_1)f$ and converge pointwise almost everywhere. \Cref{lem:density-subsolution-lone-closure} gives $\mathrm{MA}^{\mathrm{ext}}_h(u)\geq bf\,\mathrm{d}V$, attaining the maximum.
\end{proof}

The argument preserves a density lower bound. It does not claim that the original normalized solutions converge in capacity, that their limit has finite mass, or that the limiting inequality is an equality. These are additional requirements in the individual equation problem. The definition is unchanged if truncation is replaced by any increasing sequence of nonzero bounded densities tending to $f$: necessity bounds its limiting constants below by $b_h(f)$, and the same upper-envelope construction bounds them above by $b_h(f)$.

\subsection{The greatest normalized threshold subsolution}\label{subsec:threshold-envelope}

Let $X$ be a connected compact complex manifold of dimension $n\geq1$, and let $h$ be a Hermitian form on $X$. Write $\mathrm{MA}^{\mathrm{ext}}_h(u)$ for the extended non-pluripolar measure of $u\in\PSH(X,h)$, and $\mathrm{MA}_h(u)$ when $u$ is bounded. The quantities $v_-(h)$ and $v_+(h)$ are respectively the infimum and supremum of the masses of bounded $h$-plurisubharmonic potentials. Fix a smooth strictly positive probability volume form $\mathrm{d}V$; all function norms below use this volume. Fix a nonzero nonnegative $f\in L^1(\mathrm{d}V)$ and use its subsolution threshold $b_h(f)$ from \eqref{eq:individual-truncated-normalizer}.

\begin{proposition}\label{prop:greatest-normalized-threshold}
For $b=b_h(f)>0$, the family
\[
 \mathcal S_f=\{s\in\PSH(X,h):\sup_Xs=0,\
             \mathrm{MA}^{\mathrm{ext}}_h(s)\geq bf\,\mathrm{d}V\}
\]
is nonempty and compact in $L^1(X)$. It has a greatest member
\[
 \Phi_f=\left(\sup_{s\in\mathcal S_f}s\right)^*.
\]
Every normalized $L^1$ cluster limit of the greatest exponential solutions from \cref{cor:individual-positive-threshold-exponential} belongs to $\mathcal S_f$ and hence lies below $\Phi_f$.
\end{proposition}
\begin{proof}
The family is nonempty by \cref{thm:individual-density-subsolution}. Its compactness follows from normalized compactness and \cref{lem:density-subsolution-lone-closure}. Ordered contact makes it stable under finite maxima. Choose a countable Choquet family for its supremum and take successive finite maxima, including one fixed member of $\mathcal S_f$. This increasing sequence has a genuine lower bound and upper bound zero. Its canonical limit is $\Phi_f$ and its convergence holds in $L^1$. The same lemma gives the required measure inequality. Its supremum is zero because it lies between a normalized member and zero. Thus $\Phi_f\in\mathcal S_f$ and it dominates every member.

For an exponential cluster limit, the densities of the measures converge in $L^1$ to $bf$ by \cref{cor:individual-positive-threshold-exponential}. Apply the lemma to these densities and the normalized potentials.
\end{proof}

This proposition asserts a greatest normalized subsolution, not equality of its measure with $bf\,\mathrm{d}V$. In particular it does not assert that $\Phi_f$ has finite mass.


\Needspace{12\baselineskip}
\begin{proposition}\label{prop:individual-normalizer-lp}
If $f\in L^p(\mathrm{d}V)$ for some $p>1$, then $b_h(f)$ equals the normalization constant of the continuous equation for $f$. In general,
\[
 b_h(f)\geq\frac{v_-(h)}{\int_Xf\,\mathrm{d}V}.
\]
If $h$ is K\"ahler, then $b_h(f)=\int_Xh^n/\int_Xf\,\mathrm{d}V$ for every nonzero nonnegative $L^1$ density.
\end{proposition}
\begin{proof}
Let $v$ be a continuous solution for $f\in L^p$, with scalar $c>0$. It is a subsolution for $cf\,\mathrm{d}V$, so the maximal-subsolution characterization gives $c\leq b_h(f)$. Conversely, the bounded-solution comparison in \cref{lem:bounded-normalizer-subsolution}, applied to the attained threshold subsolution and $v$, gives $b_h(f)\leq c$.

Each bounded solution has total mass at least $v_-(h)$, so $c_h(f_k)\int_Xf_k\,\mathrm{d}V\geq v_-(h)$. Take the limit. In the K\"ahler case, bounded Stokes instead gives the exact identity $c_h(f_k)\int_Xf_k\,\mathrm{d}V=\int_Xh^n$, which proves the last assertion.
\end{proof}

\subsection{Smooth comparison tests}\label{subsec:individual-smooth-tests}

Bounded measurable tests in \eqref{eq:individual-bounded-tests} detect concentrations that smooth tests can miss. To compare them, let $f$ be a probability density bounded below by a positive constant and define
\[
 C_h^{\mathrm{sm}}(f)=\inf_g c_h(g)\|g/f\|_{L^\infty},
\]
where $g$ ranges over smooth strictly positive probability densities. The subsolution bound gives $b_h(f)\leq C_h^{\mathrm{sm}}(f)$, with equality for smooth $f$. The next construction separates even their positivity when lower mass vanishes.

\Needspace{10\baselineskip}
\begin{theorem}\label{thm:smooth-tests-insufficient}
If $v_-(h)=0$, there is a probability density $f\in L^1(\mathrm{d}V)$ bounded below by a positive constant such that
\[
 b_h(f)=0<C_h^{\mathrm{sm}}(f).
\]
Thus no positive multiple of $f\,\mathrm{d}V$ is dominated by an extended Monge--Amp\`ere measure, although every smooth comparison test gives a common positive bound.
\end{theorem}
\begin{proof}
Let $(U_i)_{i\geq1}$ be a countable basis of nonempty open subsets of $X$. We construct numbers $\alpha_j>0$, bounded densities $q_j\not\equiv0$, and closed nowhere dense sets $K_j$ such that
\[
 \int_Xq_j\,\mathrm{d}V\leq1,\qquad
 \operatorname{supp}q_j\subset K_j,\qquad
 c_h(q_j)\leq\alpha_j2^{-j}.
\]
Start with $\alpha_1=1/4$. Once $\alpha_j$ is chosen, the identity between lower mass and the infimum of smooth normalization constants gives a smooth strictly positive probability density $g_j$ with $c_h(g_j)<\alpha_j2^{-j}$. Choose a closed nowhere dense $K_j$ so that $q_j=g_j\mathbf1_{K_j}$ is sufficiently close to $g_j$ in $L^2$ for the displayed scalar bound to hold, using \cref{lem:lp-normalizer-continuity}. Such sets exist: remove an open dense set of arbitrarily small measure from $X$; boundedness of the fixed $g_j$ makes the $L^2$ error arbitrarily small. Choose it small enough also that $q_j$ is nonzero.

After choosing $K_1,\ldots,K_N$, put
\[
 m_N=\min_{1\leq i\leq N}
    \mathrm{d}V\left(U_i\setminus\bigcup_{j=1}^NK_j\right)>0.
\]
Each set inside the minimum is nonempty and open. Choose
\[
 0<\alpha_{N+1}\leq
       \min\{\alpha_N/2,\ 2^{-N-2}m_N\},
\]
and continue. In particular,
\[
 \sum_{j>N}\alpha_j\leq2\alpha_{N+1}
             \leq2^{-N-1}m_N.
\]
The function $Q=\sum_j\alpha_jq_j$ is nonnegative and integrable. Its essential infimum on every nonempty open set is zero. To see this on $U_i$, take $N\geq i$ and remove the first $N$ sets $K_j$. On the remaining set the first $N$ summands vanish. Markov's inequality bounds the measure where the tail exceeds $2^{-N}$ by
\[
 2^N\sum_{j>N}\alpha_j\int_Xq_j\,\mathrm{d}V
       \leq m_N/2.
\]
The remaining set originally has measure at least $m_N$. A positive-measure subset therefore satisfies $Q\leq2^{-N}$. Let $N\to\infty$ and use the basis property.

\Needspace{8\baselineskip}
Set $M=\int_X(1+Q)\,\mathrm{d}V$ and $f=(1+Q)/M$. If $\mathrm{MA}^{\mathrm{ext}}_h(u)\geq af\,\mathrm{d}V$ with $a>0$, then the bounded comparison lemma gives
\[
 \frac{a\alpha_j}{M}\leq c_h(q_j)\leq\alpha_j2^{-j}
 \quad\text{for every }j.
\]
This is impossible. Hence $b_h(f)=0$.

Finally let $g$ be any smooth strictly positive probability density and put $R=\|g/f\|_\infty$. The inequality $g\leq R(1+Q)/M$ almost everywhere forces $g\leq R/M$ everywhere. Otherwise continuity would give a nonempty open set on which $g>R(1+\varepsilon)/M$ for some $\varepsilon>0$, contradicting the essential infimum property of $Q$ on that set. Order and homogeneity of the scalar now give
\[
 R c_h(g)\geq M c_h(1).
\]
The constant test $g=1$ attains this value, since $\operatorname*{ess\,inf}_XQ=0$. Thus
$C_h^{\mathrm{sm}}(f)=M c_h(1)>0$, as required.
\end{proof}

The distinction is caused by the order restriction on a smooth test: it cannot use a measurable increase of a density whose essential lower value is unchanged in every open set. Bounded measurable tests retain this information. The theorem concerns subsolution existence for individual $L^1$ densities. A characterization of exact normalized equations at $b_h(f)$, including their natural singularity class and finite-mass compactness, remains a separate problem.
